%% HAND-WRITTEN fragment -- NOT generated by maestro2tex. Input directly from
%% AnalogChapter.tex, immediately after note_models.tex: this is the nonlinear
%% counterpart of the linear cell defined there, and both are shared by the
%% benches rather than owned by one block.
%%
%% Sources of the statements below, so they can be re-checked:
%%   model source + equation-to-reference map:
%%     ic/castalia/electrode_bc/veriloga/veriloga.va (header)
%%   grid constants (nseg, h0, beta, delta): same header, "GRID:" line
%%   wrapper topology + parameter pass-through:
%%     simruns/cv_bc_first/.../ps_cv/netlist/input.scs, subckt randles_cell_bc
%%   plan/handoff: ~/vestarv/docs/afe_rev2_cv_electrode.html
%% Parameter defaults live in Table~\ref{tab:electrode-params} -- do not
%% restate them in prose or captions.

\subsection{Voltammetric Electrode Model}\label{ss:analog-electrode-bc}

This section documents a \emph{simulation model}, not a circuit block and not
silicon: \texttt{electrode\_bc} is a Verilog-A element in the \texttt{castalia}
library that stands in for the electrode--electrolyte interface when a bench
sweeps potential and expects a voltammetric peak. Everything reported here is
either a standalone Spectre run of the model against closed-form
electrochemistry, or one transient of the model driven by the rev-2 channel.

The linear models of Section~\ref{ss:analog-models} cannot produce such a peak.
\(R_u + R_{\mathrm{ct}} \parallel C_{\mathrm{dl}}\) is a linear network, so
under a triangular sweep every branch current is a linear functional of the
excitation: a ramp through \(R_{\mathrm{ct}}\) plus a constant
\(C_{\mathrm{dl}}\,\mathrm{d}v/\mathrm{d}t\) box that changes sign at the
vertices. Changing element values moves slopes and never peak positions, so the
extrema stay pinned at the switching potentials
(Figure~\ref{fig:electrode-validation}, dashed). A peak needs two
nonlinearities acting together: Butler--Volmer kinetics, exponential in
overpotential about the couple's formal potential \(E^{0\prime}\), and
diffusional depletion, which collapses the surface concentration as current
flows and turns the exponential rise over into a \(t^{-1/2}\) decay.
Butler--Volmer alone, with the surface concentration pinned at bulk, gives a
sigmoid rather than a peak.

\texttt{electrode\_bc} implements both. Two coupled one-dimensional diffusion
ladders, one per species, share a single Butler--Volmer surface flux. Each
ladder is discretised by the box (finite-volume) method on an exponentially
expanding grid, \(h_j = h_0\beta^{\,j}\) with 28 slabs per species,
\(h_0 = \SI{0.5}{\micro\metre}\) and \(\beta = 1.280\); one slab is one internal
electrical node whose potential carries the slab concentration normalised by
\SI{1e-6}{\mole\per\cubic\centi\metre}, because raw
\si{\mole\per\cubic\centi\metre} sits below the simulator's default
\SI{1}{\micro\volt} voltage tolerance and solves to noise. The outer boundary is
held at bulk concentration at a finite Nernst-layer thickness
\(\delta = h_0(\beta^{28}-1)/(\beta-1) = \SI{1.79}{\milli\metre}\). Parameters
are listed in Table~\ref{tab:electrode-params}; they are in the CGS units of the
source literature (\si{\centi\metre\per\second},
\si{\square\centi\metre\per\second}, \si{\mole\per\cubic\centi\metre}), not SI.
The element emits anodic-positive current, which is the polarity the
transimpedance amplifier measures and the opposite of the cathodic-positive
convention of the reference text.

The benches instantiate it through the wrapper \texttt{randles\_cell\_bc}, which
is the topology of Figure~\ref{fig:randles-cell} with the charge-transfer
resistor replaced: \(R_{\mathrm{ce}}\), \(R_u\) and \(C_{\mathrm{dl}}\) remain
linear elements and \texttt{electrode\_bc} sits between
\(\mathrm{WE_s}\) and WE. The linear cell of Section~\ref{ss:analog-models}
stays the correct model for small-signal impedance work; only voltammetry needs
this one.
